\documentclass[11pt]{article}
%% ============================================================
%%  Shared preamble for the daily office orders
%%  (morning, midday, evening, night)
%%
%%  Each order file should begin with:
%%    \documentclass[11pt]{article}
%%    %% ============================================================
%%  Shared preamble for the daily office orders
%%  (morning, midday, evening, night)
%%
%%  Each order file should begin with:
%%    \documentclass[11pt]{article}
%%    %% ============================================================
%%  Shared preamble for the daily office orders
%%  (morning, midday, evening, night)
%%
%%  Each order file should begin with:
%%    \documentclass[11pt]{article}
%%    \input{worship-preamble.tex}
%%    \begin{document}
%%    ...
%%    \end{document}
%% ============================================================

%% ---------- Page: full letter (8.5 x 11 in) ----------
\usepackage[paperwidth=8.5in, paperheight=11in,
            top=0.7in, bottom=0.75in, inner=0.65in, outer=0.6in]{geometry}

%% ---------- Fonts ----------
\usepackage{fontspec}
\IfFontExistsTF{EBGaramond-Regular.otf}
  {\setmainfont{EBGaramond}[
     Extension=.otf, UprightFont=*-Regular, ItalicFont=*-Italic,
     BoldFont=*-Bold, BoldItalicFont=*-BoldItalic, Numbers=OldStyle]}
  {\setmainfont{texgyrepagella}[
     Extension=.otf, UprightFont=*-regular, ItalicFont=*-italic,
     BoldFont=*-bold, BoldItalicFont=*-bolditalic, Numbers=OldStyle]}

\usepackage{microtype}
\usepackage{xcolor}
\usepackage{fancyhdr}
\usepackage{needspace}

%% ---------- Colors ----------
%% rubric and rule are intentionally the same for now;
%% keep separate names so they can diverge later.
\definecolor{rubric}{RGB}{158,27,32}
\definecolor{rule}{RGB}{158,27,32}

%% ---------- Paragraph setup ----------
\setlength{\parindent}{0pt}
\setlength{\parskip}{0pt}
\linespread{1.08}
\raggedright

%% ---------- Running head / foot ----------
\pagestyle{fancy}
\fancyhf{}
\renewcommand{\headrulewidth}{0pt}
\fancyfoot[C]{\small\color{rubric}\thepage}
\fancypagestyle{plain}{\fancyhf{}\fancyfoot[C]{\small\color{rubric}\thepage}}

%% ============================================================
%%  Liturgical macros
%% ============================================================

% Title block
\newcommand{\OrderTitle}[1]{%
  \begin{center}
    {\LARGE\scshape\addfontfeatures{LetterSpace=6}#1\par}
    \vspace{4pt}
    {\color{rule}\rule{0.35\linewidth}{0.6pt}\par}
  \end{center}}
\newcommand{\OrderSubtitle}[1]{%
  \begin{center}\vspace{-6pt}{\itshape #1\par}\end{center}\vspace{6pt}}

% Act of worship heading
\newcommand{\Act}[1]{%
  \par\needspace{4\baselineskip}\vspace{12pt}%
  {\scshape\addfontfeatures{LetterSpace=4}\lowercase{#1}\par}%
  \vspace{3pt}}

% Rubric: red italic direction
\newcommand{\Rubric}[1]{%
  {\color{rubric}\itshape #1\par}\vspace{4pt}}

% Inline rubric (no paragraph break)
\newcommand{\rub}[1]{{\color{rubric}\itshape #1}}

% Speaker lines: leader in roman, people in bold
% NOTE: \L overwrites the standard LaTeX Ł character command.
%       Acceptable here since the office texts never need Ł.
\renewcommand{\L}[1]{\par\hspace{\psalmindent}#1\par\vspace{2pt}\setlength{\psalmindent}{0pt}}
% \P is defined by textcomp (loaded via fontspec) as the pilcrow ¶.
% We repurpose it for the "People" speaker line.
\renewcommand{\P}[1]{\par\hspace{\psalmindent}\textbf{#1}\par\vspace{2pt}\setlength{\psalmindent}{0pt}}

% Litany environment
\newenvironment{Litany}{\par\vspace{2pt}}{\par\vspace{4pt}}

% Psalm: indented verses, optional response
\newenvironment{Psalm}[1]{%
  \par\vspace{2pt}{\itshape #1\par}\vspace{3pt}%
  \begin{list}{}{\leftmargin=1.2em\itemindent=-1.2em
                 \itemsep=3pt\parsep=0pt\topsep=0pt
                 \parindent=1.2em}}
  {\end{list}\vspace{4pt}}
\newlength{\psalmindent}
\newcommand{\V}[1]{\item\hspace{\psalmindent}#1\setlength{\psalmindent}{0pt}}
\newcommand{\Resp}[1]{\item\textbf{#1}}

% Verse half-line marker
\newcommand{\pt}{\ {\color{rubric}*}\ }

% Indent the start of a line (use at the beginning of a source line)
\newcommand{\Indent}{\setlength{\psalmindent}{1.2em}}

% Hymn announcement
% Grouped so \hangindent/\hangafter don't leak into following text.
\newcommand{\Hymn}[4]{%
  {\hangindent=1.2em\hangafter=1
  \textit{#1}\hfill{\small\scshape #2}\par
  {\small\color{rubric}#3\hfill #4\par}}\vspace{4pt}}

% Silence
\newcommand{\Silence}{%
  \par\vspace{2pt}\begin{center}\rub{Silence}\end{center}\vspace{-2pt}}

% Amen in bold
\newcommand{\Amen}{\P{Amen.}}

% Scripture citation line
\newcommand{\Reading}[1]{{\itshape #1\par}\vspace{3pt}}

%%    \begin{document}
%%    ...
%%    \end{document}
%% ============================================================

%% ---------- Page: full letter (8.5 x 11 in) ----------
\usepackage[paperwidth=8.5in, paperheight=11in,
            top=0.7in, bottom=0.75in, inner=0.65in, outer=0.6in]{geometry}

%% ---------- Fonts ----------
\usepackage{fontspec}
\IfFontExistsTF{EBGaramond-Regular.otf}
  {\setmainfont{EBGaramond}[
     Extension=.otf, UprightFont=*-Regular, ItalicFont=*-Italic,
     BoldFont=*-Bold, BoldItalicFont=*-BoldItalic, Numbers=OldStyle]}
  {\setmainfont{texgyrepagella}[
     Extension=.otf, UprightFont=*-regular, ItalicFont=*-italic,
     BoldFont=*-bold, BoldItalicFont=*-bolditalic, Numbers=OldStyle]}

\usepackage{microtype}
\usepackage{xcolor}
\usepackage{fancyhdr}
\usepackage{needspace}

%% ---------- Colors ----------
%% rubric and rule are intentionally the same for now;
%% keep separate names so they can diverge later.
\definecolor{rubric}{RGB}{158,27,32}
\definecolor{rule}{RGB}{158,27,32}

%% ---------- Paragraph setup ----------
\setlength{\parindent}{0pt}
\setlength{\parskip}{0pt}
\linespread{1.08}
\raggedright

%% ---------- Running head / foot ----------
\pagestyle{fancy}
\fancyhf{}
\renewcommand{\headrulewidth}{0pt}
\fancyfoot[C]{\small\color{rubric}\thepage}
\fancypagestyle{plain}{\fancyhf{}\fancyfoot[C]{\small\color{rubric}\thepage}}

%% ============================================================
%%  Liturgical macros
%% ============================================================

% Title block
\newcommand{\OrderTitle}[1]{%
  \begin{center}
    {\LARGE\scshape\addfontfeatures{LetterSpace=6}#1\par}
    \vspace{4pt}
    {\color{rule}\rule{0.35\linewidth}{0.6pt}\par}
  \end{center}}
\newcommand{\OrderSubtitle}[1]{%
  \begin{center}\vspace{-6pt}{\itshape #1\par}\end{center}\vspace{6pt}}

% Act of worship heading
\newcommand{\Act}[1]{%
  \par\needspace{4\baselineskip}\vspace{12pt}%
  {\scshape\addfontfeatures{LetterSpace=4}\lowercase{#1}\par}%
  \vspace{3pt}}

% Rubric: red italic direction
\newcommand{\Rubric}[1]{%
  {\color{rubric}\itshape #1\par}\vspace{4pt}}

% Inline rubric (no paragraph break)
\newcommand{\rub}[1]{{\color{rubric}\itshape #1}}

% Speaker lines: leader in roman, people in bold
% NOTE: \L overwrites the standard LaTeX Ł character command.
%       Acceptable here since the office texts never need Ł.
\renewcommand{\L}[1]{\par\hspace{\psalmindent}#1\par\vspace{2pt}\setlength{\psalmindent}{0pt}}
% \P is defined by textcomp (loaded via fontspec) as the pilcrow ¶.
% We repurpose it for the "People" speaker line.
\renewcommand{\P}[1]{\par\hspace{\psalmindent}\textbf{#1}\par\vspace{2pt}\setlength{\psalmindent}{0pt}}

% Litany environment
\newenvironment{Litany}{\par\vspace{2pt}}{\par\vspace{4pt}}

% Psalm: indented verses, optional response
\newenvironment{Psalm}[1]{%
  \par\vspace{2pt}{\itshape #1\par}\vspace{3pt}%
  \begin{list}{}{\leftmargin=1.2em\itemindent=-1.2em
                 \itemsep=3pt\parsep=0pt\topsep=0pt
                 \parindent=1.2em}}
  {\end{list}\vspace{4pt}}
\newlength{\psalmindent}
\newcommand{\V}[1]{\item\hspace{\psalmindent}#1\setlength{\psalmindent}{0pt}}
\newcommand{\Resp}[1]{\item\textbf{#1}}

% Verse half-line marker
\newcommand{\pt}{\ {\color{rubric}*}\ }

% Indent the start of a line (use at the beginning of a source line)
\newcommand{\Indent}{\setlength{\psalmindent}{1.2em}}

% Hymn announcement
% Grouped so \hangindent/\hangafter don't leak into following text.
\newcommand{\Hymn}[4]{%
  {\hangindent=1.2em\hangafter=1
  \textit{#1}\hfill{\small\scshape #2}\par
  {\small\color{rubric}#3\hfill #4\par}}\vspace{4pt}}

% Silence
\newcommand{\Silence}{%
  \par\vspace{2pt}\begin{center}\rub{Silence}\end{center}\vspace{-2pt}}

% Amen in bold
\newcommand{\Amen}{\P{Amen.}}

% Scripture citation line
\newcommand{\Reading}[1]{{\itshape #1\par}\vspace{3pt}}

%%    \begin{document}
%%    ...
%%    \end{document}
%% ============================================================

%% ---------- Page: full letter (8.5 x 11 in) ----------
\usepackage[paperwidth=8.5in, paperheight=11in,
            top=0.7in, bottom=0.75in, inner=0.65in, outer=0.6in]{geometry}

%% ---------- Fonts ----------
\usepackage{fontspec}
\IfFontExistsTF{EBGaramond-Regular.otf}
  {\setmainfont{EBGaramond}[
     Extension=.otf, UprightFont=*-Regular, ItalicFont=*-Italic,
     BoldFont=*-Bold, BoldItalicFont=*-BoldItalic, Numbers=OldStyle]}
  {\setmainfont{texgyrepagella}[
     Extension=.otf, UprightFont=*-regular, ItalicFont=*-italic,
     BoldFont=*-bold, BoldItalicFont=*-bolditalic, Numbers=OldStyle]}

\usepackage{microtype}
\usepackage{xcolor}
\usepackage{fancyhdr}
\usepackage{needspace}

%% ---------- Colors ----------
%% rubric and rule are intentionally the same for now;
%% keep separate names so they can diverge later.
\definecolor{rubric}{RGB}{158,27,32}
\definecolor{rule}{RGB}{158,27,32}

%% ---------- Paragraph setup ----------
\setlength{\parindent}{0pt}
\setlength{\parskip}{0pt}
\linespread{1.08}
\raggedright

%% ---------- Running head / foot ----------
\pagestyle{fancy}
\fancyhf{}
\renewcommand{\headrulewidth}{0pt}
\fancyfoot[C]{\small\color{rubric}\thepage}
\fancypagestyle{plain}{\fancyhf{}\fancyfoot[C]{\small\color{rubric}\thepage}}

%% ============================================================
%%  Liturgical macros
%% ============================================================

% Title block
\newcommand{\OrderTitle}[1]{%
  \begin{center}
    {\LARGE\scshape\addfontfeatures{LetterSpace=6}#1\par}
    \vspace{4pt}
    {\color{rule}\rule{0.35\linewidth}{0.6pt}\par}
  \end{center}}
\newcommand{\OrderSubtitle}[1]{%
  \begin{center}\vspace{-6pt}{\itshape #1\par}\end{center}\vspace{6pt}}

% Act of worship heading
\newcommand{\Act}[1]{%
  \par\needspace{4\baselineskip}\vspace{12pt}%
  {\scshape\addfontfeatures{LetterSpace=4}\lowercase{#1}\par}%
  \vspace{3pt}}

% Rubric: red italic direction
\newcommand{\Rubric}[1]{%
  {\color{rubric}\itshape #1\par}\vspace{4pt}}

% Inline rubric (no paragraph break)
\newcommand{\rub}[1]{{\color{rubric}\itshape #1}}

% Speaker lines: leader in roman, people in bold
% NOTE: \L overwrites the standard LaTeX Ł character command.
%       Acceptable here since the office texts never need Ł.
\renewcommand{\L}[1]{\par\hspace{\psalmindent}#1\par\vspace{2pt}\setlength{\psalmindent}{0pt}}
% \P is defined by textcomp (loaded via fontspec) as the pilcrow ¶.
% We repurpose it for the "People" speaker line.
\renewcommand{\P}[1]{\par\hspace{\psalmindent}\textbf{#1}\par\vspace{2pt}\setlength{\psalmindent}{0pt}}

% Litany environment
\newenvironment{Litany}{\par\vspace{2pt}}{\par\vspace{4pt}}

% Psalm: indented verses, optional response
\newenvironment{Psalm}[1]{%
  \par\vspace{2pt}{\itshape #1\par}\vspace{3pt}%
  \begin{list}{}{\leftmargin=1.2em\itemindent=-1.2em
                 \itemsep=3pt\parsep=0pt\topsep=0pt
                 \parindent=1.2em}}
  {\end{list}\vspace{4pt}}
\newlength{\psalmindent}
\newcommand{\V}[1]{\item\hspace{\psalmindent}#1\setlength{\psalmindent}{0pt}}
\newcommand{\Resp}[1]{\item\textbf{#1}}

% Verse half-line marker
\newcommand{\pt}{\ {\color{rubric}*}\ }

% Indent the start of a line (use at the beginning of a source line)
\newcommand{\Indent}{\setlength{\psalmindent}{1.2em}}

% Hymn announcement
% Grouped so \hangindent/\hangafter don't leak into following text.
\newcommand{\Hymn}[4]{%
  {\hangindent=1.2em\hangafter=1
  \textit{#1}\hfill{\small\scshape #2}\par
  {\small\color{rubric}#3\hfill #4\par}}\vspace{4pt}}

% Silence
\newcommand{\Silence}{%
  \par\vspace{2pt}\begin{center}\rub{Silence}\end{center}\vspace{-2pt}}

% Amen in bold
\newcommand{\Amen}{\P{Amen.}}

% Scripture citation line
\newcommand{\Reading}[1]{{\itshape #1\par}\vspace{3pt}}


%% =====================================================================
%%  CANTICLE template  —  music, then call-and-response text.
%%
%%  A canticle in the UMH is a sung psalm/ode: it carries real staff
%%  notation (rendered here by ABC.js in the browser) plus leader/people
%%  or refrains. Build one by composing the primitives you already have:
%%    \Hymn    title + tune + UMH page
%%    \Melody  staff notation (new — ABC.js renders it in the browser)
%%    \L / \P  leader / people   (or a full Litany / Psalm environment)
%%
%%  To add a new canticle, copy this page, rename it, and replace the
%%  bracketed text + the TEST melody below with the real tune.
%%
%%  NOTE on the \Melody body: this is ABC notation, one or more "parts"
%%  separated by |...|| at the end.  Keep it on its own line(s); the
%%  braces in this file match, so it is safe to wrap across lines.
%%
%%  Build (HTML):  python3 tex2html.py canticle.tex
%% =====================================================================

\begin{document}

\OrderTitle{[Canticle Title]}
\OrderSubtitle{[a subtitle or occasion, if any]}

%% ---- optional rubric, e.g. "The people may stand or sit as able." ----
%% \Rubric{[direction]}

\Act{Canticle}

%% --- Hymn announcement: title, UMH page; tune name + meter optional ------
\Hymn{[Canticle title]}{UMH [nnn]}{[Tune name]}{[Meter, e.g. LM]}

%% ---- MUSIC -------------------------------------------------------------
%% \Melody holds ABC notation.  Replace this TEST scale with the canticle's
%% actual tune (see abcjs.example.com for a live editor; the notation can be
%% pasted in straight from there).  Anything that isn't a renderable ABC
%% line will be shown as source text rather than drawn — never a blank gap.
\Melody{
X:1
T:[Canticle]
M:4/4
L:1/4
K:C
C D E F | G A B c | C D E F | c4 ||
}

%% ---- CALL AND RESPONSE -------------------------------------------------
%% Any of these already work in your pipeline — pick whichever the canticle
%% uses.  Leader (roman) / People (bold) is the most common shape.

\Rubric{[Leader reads / sings a part; the people respond.] }
\begin{Litany}
\L{[Leader's line]}
\P{[People's response]}
\L{[Leader's line]}
\P{[People's response]}
\end{Litany}

%% ---- (ALTERNATE) verse-with-refrain canticle --------------------------
%% Use this instead of the Litany block for a through-sung canticle with a
%% repeating refrain between verses:
%%
%% \begin{Psalm}{[Ref., e.g. "Nunc Dimittis, Luke 2:29-32"]}
%% \Resp{[refrain or antiphon]}
%% \V{[verse line]}
%% \V{[verse line]}
%% \Resp{[refrain or antiphon]}
%% \V{[verse line]}
%% \end{Psalm}

\Silence
\Amen

\end{document}
